\documentclass[11pt,a4paper]{article}

% --- Packages ---
\usepackage[utf8]{inputenc}
\usepackage[T1]{fontenc}
\usepackage{amsmath, amssymb, amsfonts}
\usepackage{geometry}
\usepackage{hyperref}
\usepackage{listings}
\usepackage{xcolor}
\usepackage{booktabs}
\usepackage{enumitem}
\usepackage{caption}
\usepackage{float}
\usepackage{titlesec}

\geometry{margin=1in}
\hypersetup{colorlinks=true, linkcolor=blue!70!black, urlcolor=blue!70!black, citecolor=blue!70!black}

% --- Code style ---
\definecolor{codebg}{HTML}{F7F7F7}
\definecolor{codeframe}{HTML}{CCCCCC}
\definecolor{kwcolor}{HTML}{0000CC}
\definecolor{strcolor}{HTML}{008800}
\definecolor{commentcolor}{HTML}{888888}
\lstdefinestyle{pythonstyle}{
  backgroundcolor=\color{codebg},
  frame=single,
  rulecolor=\color{codeframe},
  basicstyle=\ttfamily\footnotesize,
  keywordstyle=\color{kwcolor}\bfseries,
  stringstyle=\color{strcolor},
  commentstyle=\color{commentcolor}\itshape,
  language=Python,
  breaklines=true,
  showstringspaces=false,
  tabsize=4,
  numbers=left,
  numberstyle=\tiny\color{gray},
  xleftmargin=1.5em,
  framexleftmargin=1em,
  columns=fullflexible,
}
\lstset{style=pythonstyle}

% --- Title ---
\title{\textbf{OWCM Library: Environment Submodule Reference}\\[0.3em]
\large Mathematical Formulations \& Code Snippets}
\author{Auto-generated from \texttt{owcm} source code}
\date{\today}

\begin{document}
\maketitle
\tableofcontents
\newpage

%=============================================================================
\section{Overview}
%=============================================================================

The \texttt{owcm} library organises optical wireless channel models into five
\emph{environment submodules}:

\begin{enumerate}[leftmargin=2em]
  \item \textbf{Atmospheric} --- fog attenuation, geometric loss, pointing error, turbulence
  \item \textbf{Indoor} --- Lambertian LoS/NLoS, blockage
  \item \textbf{Space} --- free-space path loss, deep-space link budget
  \item \textbf{Underground} --- dust extinction, tunnel LoS, blockage
  \item \textbf{Underwater} --- extinction, scattering, reflective/retroreflective links
\end{enumerate}

Each model class exposes three methods:
\begin{itemize}
  \item \texttt{gain(link, env)} --- deterministic channel gain $H$ (linear scale)
  \item \texttt{sample(link, env, n)} --- draw $n$ random samples from the channel
  \item \texttt{describe()} --- metadata dictionary
\end{itemize}

%=============================================================================
%=============================================================================
\newpage
\section{Atmospheric Submodule}
%=============================================================================
%=============================================================================

%-----------------------------------------------------------------------------
\subsection{Attenuation Models}
%-----------------------------------------------------------------------------

% ---- 1. Beer-Lambert ----
\subsubsection{Beer--Lambert Attenuation}
\label{sec:beer_lambert}

\paragraph{Registry key.} \texttt{("atmospheric", "attenuation", "beer\_lambert")}

\paragraph{Mathematical form.}
\begin{equation}
\boxed{H_{\text{att}} = e^{-\alpha \, R}}
\end{equation}
where $\alpha$ [\si{m^{-1}}] is the extinction coefficient and $R$ [m] is the
link range.

\paragraph{Code snippet.}
\begin{lstlisting}
from owcm.core.types import Tx, Rx, Link, Environment
from owcm.atmospheric.models.attenuation import BeerLambertAttenuation

tx = Tx(wavelength_nm=1550)
rx = Rx()
link = Link(tx=tx, rx=rx)        # range_m computed from positions
env  = Environment()

model = BeerLambertAttenuation(link, env, alpha=0.25)
H = model.gain(link, env)        # deterministic gain
print(f"Beer-Lambert gain = {H:.6f}")
\end{lstlisting}

% ---- 2. Kim ----
\subsubsection{Kim Fog Model}
\label{sec:kim}

\paragraph{Registry key.} \texttt{("atmospheric", "attenuation", "kim")}

\paragraph{Mathematical form.}
The specific attenuation in dB/km is:
\begin{equation}
A = \frac{3.91}{V} \left(\frac{\lambda}{550}\right)^{-q(V)}
\quad \text{[dB/km]}
\end{equation}
where $V$ is visibility in km, $\lambda$ is wavelength in nm, and $q(V)$ is
the Kim size-parameter:
\begin{equation}
q(V) =
\begin{cases}
1.6     & V > 50\text{ km} \\
1.3     & 6 < V \le 50\text{ km} \\
0.16V + 0.34 & 1 < V \le 6\text{ km} \\
V - 0.5 & 0.5 < V \le 1\text{ km} \\
0       & V \le 0.5\text{ km}
\end{cases}
\end{equation}
The channel gain is then:
\begin{equation}
\boxed{H_{\text{Kim}} = \exp\!\left(-\frac{A}{1000}\,R\right)}
\end{equation}
with $A$ converted from dB/km to Np/m internally.

\paragraph{Code snippet.}
\begin{lstlisting}
from owcm.atmospheric.models.attenuation import KimFogModel

model = KimFogModel(link, env, visibility_km=2.0)
H = model.gain(link, env)
print(f"Kim fog gain = {H:.6f}")
\end{lstlisting}

% ---- 3. Kruse ----
\subsubsection{Kruse Fog Model}
\label{sec:kruse}

\paragraph{Registry key.} \texttt{("atmospheric", "attenuation", "kruse")}

\paragraph{Mathematical form.}
Same attenuation expression as Kim but with a simplified $q$ exponent:
\begin{equation}
q(V) =
\begin{cases}
1.6                 & V > 50\text{ km} \\
1.3                 & 6 < V \le 50\text{ km} \\
0.585\,V^{1/3}      & V \le 6\text{ km}
\end{cases}
\end{equation}
\begin{equation}
\boxed{H_{\text{Kruse}} = \exp\!\left(-\frac{3.91}{1000\,V}\left(\frac{\lambda}{550}\right)^{-q} R\right)}
\end{equation}

\paragraph{Code snippet.}
\begin{lstlisting}
from owcm.atmospheric.models.attenuation import KruseFogModel

model = KruseFogModel(link, env, visibility_km=4.0)
H = model.gain(link, env)
print(f"Kruse fog gain = {H:.6f}")
\end{lstlisting}

% ---- 4. Al Naboulsi Convection ----
\subsubsection{Al Naboulsi Convection Fog}
\label{sec:al_naboulsi_conv}

\paragraph{Registry key.} \texttt{("atmospheric", "attenuation", "al\_naboulsi\_convection")}

\paragraph{Mathematical form.}
\begin{equation}
\boxed{A = \frac{10}{\ln 10}\;\frac{0.11478\,\lambda + 3.8367}{V}
\quad \text{[dB/km]}}
\end{equation}
where $\lambda$ is in \si{\micro m} and $V$ is visibility in km.
Channel gain from specific attenuation:
\begin{equation}
H = \exp\!\left(-\frac{A}{4.343 \times 1000}\,R\right)
\end{equation}

\paragraph{Code snippet.}
\begin{lstlisting}
from owcm.atmospheric.models.attenuation import AlNaboulsiConvection

model = AlNaboulsiConvection(link, env, visibility_km=1.0)
H = model.gain(link, env)
print(f"Al Naboulsi (convection) gain = {H:.6f}")
\end{lstlisting}

% ---- 5. Al Naboulsi Advection ----
\subsubsection{Al Naboulsi Advection Fog}
\label{sec:al_naboulsi_adv}

\paragraph{Registry key.} \texttt{("atmospheric", "attenuation", "al\_naboulsi\_advection")}

\paragraph{Mathematical form.}
\begin{equation}
\boxed{A = \frac{10}{\ln 10}\;\frac{0.18126\,\lambda^2 + 0.13709\,\lambda + 3.7502}{V}
\quad \text{[dB/km]}}
\end{equation}
with identical gain conversion as the convection variant.

\paragraph{Code snippet.}
\begin{lstlisting}
from owcm.atmospheric.models.attenuation import AlNaboulsiAdvection

model = AlNaboulsiAdvection(link, env, visibility_km=0.8)
H = model.gain(link, env)
print(f"Al Naboulsi (advection) gain = {H:.6f}")
\end{lstlisting}

% ---- 6. Ijaz ----
\subsubsection{Ijaz Controlled Fog}
\label{sec:ijaz}

\paragraph{Registry key.} \texttt{("atmospheric", "attenuation", "ijaz")}

\paragraph{Mathematical form.}
\begin{equation}
q(\lambda) = 0.1428\,\lambda - 0.0947
\end{equation}
\begin{equation}
\boxed{A = \frac{17}{V}\left(\frac{\lambda}{0.55}\right)^{-q(\lambda)}
\quad \text{[dB/km]}}
\end{equation}
where $\lambda$ is in \si{\micro m}.

\paragraph{Code snippet.}
\begin{lstlisting}
from owcm.atmospheric.models.attenuation import IjazFog

model = IjazFog(link, env, visibility_km=0.5)
H = model.gain(link, env)
print(f"Ijaz fog gain = {H:.6f}")
\end{lstlisting}

% ---- 7. Esmail ----
\subsubsection{Esmail Unified Fog Regression}
\label{sec:esmail}

\paragraph{Registry key.} \texttt{("atmospheric", "attenuation", "esmail")}

\paragraph{Mathematical form.}
\begin{equation}
\boxed{A = k\,V^{-(a\lambda + b)} \quad \text{[dB/km]}}
\end{equation}
The regression constants $k$, $a$, $b$ must be supplied by the user (they are
\emph{not} hardcoded).

\paragraph{Code snippet.}
\begin{lstlisting}
from owcm.atmospheric.models.attenuation import EsmailFog

model = EsmailFog(link, env,
                  visibility_km=1.0,
                  k=14.0, a=0.1, b=0.6)
H = model.gain(link, env)
print(f"Esmail fog gain = {H:.6f}")
\end{lstlisting}

%-----------------------------------------------------------------------------
\subsection{Geometric Loss Models}
%-----------------------------------------------------------------------------

% ---- 8. Uniform Aperture ----
\subsubsection{Uniform Aperture Geometric Loss}
\label{sec:uniform_aperture}

\paragraph{Registry key.} \texttt{("atmospheric", "geometry", "uniform\_aperture")}

\paragraph{Mathematical form.}
\begin{equation}
D_B = D_T + R\,\theta
\end{equation}
\begin{equation}
\boxed{L_{\text{geom}} = \left(\frac{D_R}{D_B}\right)^2}
\end{equation}
where $D_T,\,D_R$ are the transmitter and receiver aperture diameters, $R$ is
range, and $\theta$ is the beam divergence angle.

\paragraph{Code snippet.}
\begin{lstlisting}
from owcm.atmospheric.models.geometry import UniformApertureGeometricLoss

tx = Tx(aperture_diameter_m=0.05, divergence_angle_rad=1e-3)
rx = Rx(aperture_diameter_m=0.1)
link = Link(tx=tx, rx=rx)
model = UniformApertureGeometricLoss(link, env)
H = model.gain(link, env)
print(f"Uniform aperture geom. loss = {H:.6e}")
\end{lstlisting}

% ---- 9. Gaussian Exact ----
\subsubsection{Gaussian Collected Fraction --- Exact}
\label{sec:gauss_exact}

\paragraph{Registry key.} \texttt{("atmospheric", "geometry", "gaussian\_collected\_exact")}

\paragraph{Mathematical form.}
\begin{equation}
w_z = \frac{D_T}{2} + R\,\theta, \quad a = \frac{D_R}{2}
\end{equation}
\begin{equation}
\boxed{A_0 = \left[\operatorname{erf}\!\left(\frac{\sqrt{2}\,a}{w_z}\right)\right]^2}
\end{equation}

\paragraph{Code snippet.}
\begin{lstlisting}
from owcm.atmospheric.models.geometry import GaussianCollectedFractionExact

model = GaussianCollectedFractionExact(link, env)
A0 = model.gain(link, env)
print(f"Gaussian exact A0 = {A0:.6f}")
\end{lstlisting}

% ---- 10. Gaussian Approx ----
\subsubsection{Gaussian Collected Fraction --- Approximate}
\label{sec:gauss_approx}

\paragraph{Registry key.} \texttt{("atmospheric", "geometry", "gaussian\_collected\_approx")}

\paragraph{Mathematical form.}
\begin{equation}
\boxed{A_0 \approx 1 - \exp\!\left(-\frac{2a^2}{w_z^2}\right)}
\end{equation}

\paragraph{Code snippet.}
\begin{lstlisting}
from owcm.atmospheric.models.geometry import GaussianCollectedFractionApprox

model = GaussianCollectedFractionApprox(link, env)
A0 = model.gain(link, env)
print(f"Gaussian approx A0 = {A0:.6f}")
\end{lstlisting}

%-----------------------------------------------------------------------------
\subsection{Pointing Error Models}
%-----------------------------------------------------------------------------

% ---- 11. Rayleigh-Gaussian ----
\subsubsection{Rayleigh--Gaussian Pointing Loss}
\label{sec:rayleigh_gaussian}

\paragraph{Registry key.} \texttt{("atmospheric", "pointing", "rayleigh\_gaussian")}

\paragraph{Mathematical form.}
The radial jitter $r$ follows a Rayleigh distribution with parameter
$\sigma_s$.  The instantaneous pointing loss is:
\begin{equation}
L_p(r) = \exp\!\left(-\frac{2r^2}{w_{\text{eq}}^2}\right)
\end{equation}
The \textbf{mean} pointing gain (returned by \texttt{gain()}) is:
\begin{equation}
\boxed{\langle L_p \rangle = \frac{w_{\text{eq}}^2}{w_{\text{eq}}^2 + 4\sigma_s^2}}
\end{equation}
The \texttt{sample()} method draws $r \sim \mathrm{Rayleigh}(\sigma_s)$ and
returns the per-sample loss.

\paragraph{Code snippet.}
\begin{lstlisting}
from owcm.atmospheric.models.pointing import RayleighGaussianPointingLoss

model = RayleighGaussianPointingLoss(
    link, env,
    jitter_sigma_m=0.01,
    wz_eq_m=0.1
)
mean_gain = model.gain(link, env)
samples   = model.sample(link, env, n=10000)
print(f"Mean pointing gain = {mean_gain:.4f}")
\end{lstlisting}

%-----------------------------------------------------------------------------
\subsection{Turbulence Models}
%-----------------------------------------------------------------------------

% ---- 12. Log-Normal ----
\subsubsection{Log-Normal Turbulence}
\label{sec:lognormal}

\paragraph{Registry key.} \texttt{("atmospheric", "turbulence", "lognormal")}

\paragraph{Mathematical form.}
Under weak turbulence, the irradiance $I$ follows a log-normal distribution:
\begin{equation}
\ln I \sim \mathcal{N}\!\left(-\tfrac{1}{2}\sigma_R^2,\;\sigma_R^2\right)
\end{equation}
so that $\mathbb{E}[I]=1$.
\begin{equation}
\boxed{h_t = \exp\!\left(\mathcal{N}\!\left(-\tfrac{\sigma_R^2}{2},\;\sigma_R^2\right)\right)}
\end{equation}
\texttt{gain()} returns 1 (mean); \texttt{sample()} draws random fading
realisations.

\paragraph{Code snippet.}
\begin{lstlisting}
from owcm.atmospheric.models.turbulence import LogNormalTurbulence

model = LogNormalTurbulence(link, env, sigma_R2=0.3)
samples = model.sample(link, env, n=10000)
print(f"Mean of samples = {samples.mean():.4f}")
\end{lstlisting}

% ---- 13. Gamma-Gamma ----
\subsubsection{Gamma--Gamma Turbulence}
\label{sec:gamma_gamma}

\paragraph{Registry key.} \texttt{("atmospheric", "turbulence", "gamma\_gamma")}

\paragraph{Mathematical form.}
For moderate-to-strong turbulence, the irradiance is the product of two
independent Gamma random variables:
\begin{equation}
I = X \cdot Y, \quad
X \sim \Gamma(\alpha,\,1/\alpha), \quad
Y \sim \Gamma(\beta,\,1/\beta)
\end{equation}
\begin{equation}
\boxed{f_I(I) = \frac{2(\alpha\beta)^{(\alpha+\beta)/2}}
{\Gamma(\alpha)\,\Gamma(\beta)}\;
I^{(\alpha+\beta)/2-1}\;
K_{\alpha-\beta}\!\left(2\sqrt{\alpha\beta\,I}\right)}
\end{equation}

\paragraph{Code snippet.}
\begin{lstlisting}
from owcm.atmospheric.models.turbulence import GammaGammaTurbulence

model = GammaGammaTurbulence(link, env, gg_alpha=4.0, gg_beta=1.5)
samples = model.sample(link, env, n=10000)
print(f"Gamma-Gamma mean = {samples.mean():.4f}")
\end{lstlisting}


%=============================================================================
%=============================================================================
\newpage
\section{Indoor Submodule}
%=============================================================================
%=============================================================================

%-----------------------------------------------------------------------------
\subsection{LoS Models}
%-----------------------------------------------------------------------------

% ---- 14. Lambertian LoS ----
\subsubsection{Lambertian LoS Channel Gain}
\label{sec:indoor_lambertian_los}

\paragraph{Registry key.} \texttt{("indoor", "los", "lambertian\_los")}

\paragraph{Mathematical form.}
(FSO\_Survey Eq.\,27--31.)  The Lambertian order is:
\begin{equation}
m = -\frac{\ln 2}{\ln(\cos \phi_{1/2})}
\end{equation}
The radiation intensity pattern is:
\begin{equation}
R_o(\phi) = \frac{m+1}{2\pi}\,\cos^m(\phi)
\end{equation}
The optical concentrator gain is:
\begin{equation}
g(\omega) = \frac{n^2}{\sin^2(\omega_c)}, \quad 0 \le \omega \le \omega_c
\end{equation}
The LoS DC gain is:
\begin{equation}
\boxed{H_{\text{LoS}} = \frac{A_d}{d^2}\;R_o(\phi_0)\;\cos(\omega_0)\;T_s(\omega_0)\;g(\omega_0),
\quad \omega_0 \le \omega_c}
\end{equation}
where $A_d$ is the physical detector area, $d$ is the Tx--Rx distance,
$\phi_0$ is the irradiance angle, $\omega_0$ is the incidence angle,
$\omega_c$ is the receiver FOV half-angle, $T_s$ is the optical filter gain,
and $n$ is the refractive index of the concentrator.

\paragraph{Code snippet.}
\begin{lstlisting}
from owcm.indoor.models.all_models import IndoorLambertianLoS

tx  = Tx(position=np.array([2.5, 2.5, 3.0]),
         wavelength_nm=850,
         aperture_diameter_m=0.01)
rx  = Rx(position=np.array([1.0, 1.0, 0.85]),
         aperture_diameter_m=0.01,
         field_of_view_rad=np.deg2rad(60))
link = Link(tx=tx, rx=rx)
env  = Environment(params={"cos_phi": 0.95, "cos_psi": 0.90})

model = IndoorLambertianLoS(link, env,
    semi_angle_half_power_rad=np.deg2rad(60),
    concentrator_refractive_index=1.5,
    optical_filter_gain=1.0)
H = model.gain(link, env)
print(f"Indoor LoS gain = {H:.6e}")
\end{lstlisting}

% ---- 15. Lambertian LoS+NLoS ----
\subsubsection{Lambertian LoS + NLoS Combined}
\label{sec:indoor_los_nlos}

\paragraph{Registry key.} \texttt{("indoor", "los", "lambertian\_los\_plus\_nlos")}

\paragraph{Mathematical form.}
(FSO\_Survey Eq.\,37.)
\begin{equation}
\boxed{H = H_{\text{LoS}} + H_{\text{NLoS}}}
\end{equation}
Internally instantiates both \texttt{IndoorLambertianLoS} and
\texttt{KBounceNLoS} and sums their gains.

\paragraph{Code snippet.}
\begin{lstlisting}
from owcm.indoor.models.all_models import IndoorLambertianLoSPlusNLoS

model = IndoorLambertianLoSPlusNLoS(link, env,
    semi_angle_half_power_rad=np.deg2rad(60),
    reflectance=0.7, room_area_m2=25.0, n_bounces=3)
H = model.gain(link, env)
print(f"LoS + NLoS combined gain = {H:.6e}")
\end{lstlisting}

%-----------------------------------------------------------------------------
\subsection{NLoS Models}
%-----------------------------------------------------------------------------

% ---- 16. K-Bounce ----
\subsubsection{K-Bounce Diffuse NLoS}
\label{sec:kbounce}

\paragraph{Registry key.} \texttt{("indoor", "nlos", "kbounce")}

\paragraph{Mathematical form.}
(FSO\_Survey Eq.\,32--36.)  The single-bounce differential gain from a
reflector patch $\Delta A$ is:
\begin{equation}
\mathrm{d}H_1 = \frac{A_d}{d_{S_k}^2 \, d_{R_k}^2}\;
R_o(\phi_{S_k})\;\rho\;\Delta A\;
\cos(\omega_{S_k})\;\cos(\phi_{R_k})\;\cos(\omega_{R_k})\;T_s\;g
\end{equation}
For $K$ bounces, gains are accumulated with a geometric decay factor $\rho$:
\begin{equation}
\boxed{H_{\text{NLoS}} = \sum_{k=1}^{K} H_1 \cdot \rho^{k-1}}
\end{equation}

\paragraph{Code snippet.}
\begin{lstlisting}
from owcm.indoor.models.all_models import KBounceNLoS

model = KBounceNLoS(link, env,
    reflectance=0.7,
    room_area_m2=25.0,
    n_bounces=3,
    ceiling_height_m=3.0)
H = model.gain(link, env)
print(f"K-bounce NLoS gain = {H:.6e}")
\end{lstlisting}

%-----------------------------------------------------------------------------
\subsection{Blockage Models}
%-----------------------------------------------------------------------------

% ---- 17. On-Off Blockage ----
\subsubsection{On--Off Blockage}
\label{sec:onoff_blockage}

\paragraph{Registry key.} \texttt{("indoor", "blockage", "onoff")}

\paragraph{Mathematical form.}
The link is either fully blocked or unblocked:
\begin{equation}
H_b \sim \text{Bernoulli}(1 - p_{\text{blocked}})
\end{equation}
\begin{equation}
\boxed{\langle H_b \rangle = 1 - p_{\text{blocked}}}
\end{equation}

\paragraph{Code snippet.}
\begin{lstlisting}
from owcm.indoor.models.all_models import OnOffBlockage

model = OnOffBlockage(link, env, p_blocked=0.15)
mean_gain = model.gain(link, env)      # 0.85
samples   = model.sample(link, env, n=1000)
\end{lstlisting}

% ---- 18. Markov Blockage ----
\subsubsection{Markov Chain Blockage}
\label{sec:markov_blockage}

\paragraph{Registry key.} \texttt{("indoor", "blockage", "markov")}

\paragraph{Mathematical form.}
Two-state Markov chain with transition probabilities:
\begin{equation}
p_{01} = P(\text{unblocked} \to \text{blocked}), \quad
p_{10} = P(\text{blocked} \to \text{unblocked})
\end{equation}
Steady-state unblocked probability:
\begin{equation}
\boxed{\pi_{\text{unblocked}} = \frac{p_{10}}{p_{01} + p_{10}}}
\end{equation}

\paragraph{Code snippet.}
\begin{lstlisting}
from owcm.indoor.models.all_models import MarkovBlockage

model = MarkovBlockage(link, env, p_01=0.1, p_10=0.3)
ss_gain = model.gain(link, env)       # 0.75
chain   = model.sample(link, env, n=500)
\end{lstlisting}


%=============================================================================
%=============================================================================
\newpage
\section{Space Submodule}
%=============================================================================
%=============================================================================

%-----------------------------------------------------------------------------
\subsection{Free-Space Path Loss}
%-----------------------------------------------------------------------------

% ---- 19. FSPL ----
\subsubsection{Free-Space Path Loss (FSPL)}
\label{sec:fspl}

\paragraph{Registry key.} \texttt{("space", "free\_space", "fspl")}

\paragraph{Mathematical form.}
\begin{equation}
\boxed{L_{\text{fs}} = \left(\frac{\lambda}{4\pi R}\right)^2}
\end{equation}
where $\lambda$ is wavelength in metres and $R$ is range in metres.

\paragraph{Code snippet.}
\begin{lstlisting}
from owcm.space.models.free_space import FreeSpacePathLoss

tx = Tx(wavelength_nm=1550,
        position=np.array([0, 0, 0]))
rx = Rx(position=np.array([0, 0, 400e3]))
link = Link(tx=tx, rx=rx)
model = FreeSpacePathLoss(link, env)
L = model.gain(link, env)
print(f"FSPL = {L:.6e}")
\end{lstlisting}

%-----------------------------------------------------------------------------
\subsection{Deep-Space Link Budget}
%-----------------------------------------------------------------------------

% ---- 20. Table XIII ----
\subsubsection{Deep-Space Table XIII Model}
\label{sec:deepspace}

\paragraph{Registry key.} \texttt{("space", "deepspace", "table\_xiii")}

\paragraph{Mathematical form.}
(Survey Table XIII, Eq.\,112--116.)
\begin{align}
L_{\text{fs}} &= \left(\frac{\lambda}{4\pi R}\right)^2
  & &\text{(Eq.\,113)} \\[4pt]
G_t &= \frac{8}{w_0^2}
  & &\text{(Eq.\,114)} \\[4pt]
G_r &= \left(\frac{\pi D_R}{\lambda}\right)^2\!(1 - \gamma_r^2)
  & &\text{(Eq.\,115)} \\[4pt]
L_{\text{pt}} &= \left(\frac{w_0}{4\sigma_p}\right)^2 p_0
  & &\text{(Eq.\,116)}
\end{align}
The overall channel gain is:
\begin{equation}
\boxed{H = G_t \cdot G_r \cdot L_{\text{fs}} \cdot L_a \cdot L_s \cdot L_b
  \cdot L_{\text{pt}} \cdot \eta_t \cdot \eta_r}
\end{equation}

\paragraph{Code snippet.}
\begin{lstlisting}
from owcm.space.models.deepspace import DeepSpaceTableXIII

tx = Tx(wavelength_nm=1064,
        position=np.array([0, 0, 0]))
rx = Rx(aperture_diameter_m=1.0,
        position=np.array([0, 0, 384.4e6]))  # Moon
link = Link(tx=tx, rx=rx)
env = Environment()

model = DeepSpaceTableXIII(link, env,
    w0_rad=15e-6,       # beam divergence
    gamma_r=0.1,        # Rx obscuration ratio
    sigma_p=2e-6,       # pointing jitter
    p0=0.99, La=0.9, Ls=0.95, Lb=0.98,
    eta_t=0.9, eta_r=0.85)
H = model.gain(link, env)
print(f"Deep-space gain H = {H:.6e}")
\end{lstlisting}


%=============================================================================
%=============================================================================
\newpage
\section{Underground Submodule}
%=============================================================================
%=============================================================================

%-----------------------------------------------------------------------------
\subsection{LoS Model}
%-----------------------------------------------------------------------------

% ---- 21. Underground LoS ----
\subsubsection{Underground Lambertian LoS}
\label{sec:underground_los}

\paragraph{Registry key.} \texttt{("underground", "los", "lambertian\_los")}

\paragraph{Mathematical form.}
Simplified inverse-square-law decay for tunnel VLC (Survey Eq.\,57 simplified):
\begin{equation}
\boxed{H_{\text{LoS}} = \frac{1}{R^2}}
\end{equation}

\paragraph{Code snippet.}
\begin{lstlisting}
from owcm.underground.models.all_models import UndergroundLoS

model = UndergroundLoS(link, env)
H = model.gain(link, env)
print(f"Underground LoS gain = {H:.6e}")
\end{lstlisting}

%-----------------------------------------------------------------------------
\subsection{Dust Extinction}
%-----------------------------------------------------------------------------

% ---- 22. Dust Extinction ----
\subsubsection{Exponential Dust Extinction}
\label{sec:dust_extinction}

\paragraph{Registry key.} \texttt{("underground", "dust", "exp\_extinction")}

\paragraph{Mathematical form.}
\begin{equation}
\boxed{H_{\text{dust}} = \exp(-k\,C\,R)}
\end{equation}
where $C$ is the dust concentration [\si{mg/m^3}], $k$ is the specific
extinction coefficient [\si{m^{-1}/(mg/m^3)}], and $R$ is range.

\paragraph{Code snippet.}
\begin{lstlisting}
from owcm.underground.models.all_models import DustExtinction

model = DustExtinction(link, env,
    dust_conc_mg_m3=5.0,
    dust_k_per_m_per_mg_m3=0.5)
H = model.gain(link, env)
print(f"Dust extinction gain = {H:.6f}")
\end{lstlisting}

%-----------------------------------------------------------------------------
\subsection{Blockage Models}
%-----------------------------------------------------------------------------

\subsubsection{On--Off Blockage (Underground)}
\label{sec:ug_onoff}

\paragraph{Registry key.} \texttt{("underground", "blockage", "onoff")}

Identical formulation to the indoor on--off blockage (Sec.\,\ref{sec:onoff_blockage}).

\begin{lstlisting}
from owcm.underground.models.all_models import OnOffBlockage

model = OnOffBlockage(link, env, p_blocked=0.2)
H = model.gain(link, env)   # 0.8
\end{lstlisting}

\subsubsection{Markov Blockage (Underground)}
\label{sec:ug_markov}

\paragraph{Registry key.} \texttt{("underground", "blockage", "markov")}

Same two-state Markov chain as indoor (Sec.\,\ref{sec:markov_blockage}),
adapted for tunnel obstructions (machinery, vehicles).

\begin{lstlisting}
from owcm.underground.models.all_models import MarkovBlockage

model = MarkovBlockage(link, env, p_01=0.15, p_10=0.4)
H = model.gain(link, env)
\end{lstlisting}


%=============================================================================
%=============================================================================
\newpage
\section{Underwater Submodule}
%=============================================================================
%=============================================================================

%-----------------------------------------------------------------------------
\subsection{Extinction Models}
%-----------------------------------------------------------------------------

% ---- 25. a+b Extinction ----
\subsubsection{Absorption + Scattering Extinction}
\label{sec:uw_extinction}

\paragraph{Registry key.} \texttt{("underwater", "extinction", "a\_plus\_b")}

\paragraph{Mathematical form.}
The total extinction coefficient is the sum of absorption and scattering:
\begin{equation}
c(\lambda) = a(\lambda) + b(\lambda)
\end{equation}
\begin{equation}
\boxed{H_{\text{ext}} = \exp\!\left(-c\,R\right)}
\end{equation}
Default $c = 0.05\;\si{m^{-1}}$ (clear ocean).  Users may supply
\texttt{attenuation\_c} directly or both \texttt{attenuation\_a} and
\texttt{attenuation\_b} via \texttt{env.params}.

\paragraph{Code snippet.}
\begin{lstlisting}
from owcm.underwater.models.all_models import UnderwaterExtinction

env = Environment(params={"attenuation_a": 0.02, "attenuation_b": 0.03})
model = UnderwaterExtinction(link, env)
H = model.gain(link, env)
print(f"Underwater extinction gain = {H:.6f}")
\end{lstlisting}

%-----------------------------------------------------------------------------
\subsection{Survey Link Models}
%-----------------------------------------------------------------------------

% ---- 26. Table XI LoS ----
\subsubsection{Table XI LoS (Eq.\,89)}
\label{sec:uw_los}

\paragraph{Registry key.} \texttt{("underwater", "link", "los\_case")}

\paragraph{Mathematical form.}
\begin{equation}
z = \frac{d}{\cos\theta}
\end{equation}
\begin{equation}
\boxed{P_R = P_T \,\eta_T\,\eta_R \;
\frac{A_{\text{rec}}\,\cos\theta}{2\pi d^2 (1-\cos\theta_0)}
\; e^{-c\,z}}
\end{equation}
The gain is $H = P_R / P_T$.

\paragraph{Code snippet.}
\begin{lstlisting}
from owcm.underwater.models.survey_model import TableXILoS

model = TableXILoS(link, env,
    theta_rad=0.0,
    theta0_rad=np.deg2rad(10),
    etaT=0.9, etaR=0.9,
    attenuation_c=0.05)
H = model.gain(link, env)
print(f"Table XI LoS gain = {H:.6e}")
\end{lstlisting}

% ---- 27. Table XI Retroreflector ----
\subsubsection{Table XI Retroreflector (Eq.\,90)}
\label{sec:uw_retro}

\paragraph{Registry key.} \texttt{("underwater", "link", "retroreflective\_case")}

\paragraph{Mathematical form.}
The round-trip propagation distance is $z = 2d/\cos\theta$.
\begin{equation}
\boxed{P_R = P_T\,\eta_T\,\eta_R\,\eta_{\text{retro}}\;
e^{-c\,z}\;
\frac{A_{\text{retro}}\cos\theta}
     {2\pi d^2 (1-\cos\theta_0)}
\;\frac{A_{\text{rec}}\cos\theta}
       {\pi (d\tan\theta_{0,r})^2}}
\end{equation}

\paragraph{Code snippet.}
\begin{lstlisting}
from owcm.underwater.models.survey_model import TableXIRetro

model = TableXIRetro(link, env,
    theta_rad=0.0,
    theta0_rad=np.deg2rad(10),
    theta0_retro_rad=np.deg2rad(0.5),
    etaT=0.9, etaR=0.9, etaRetro=0.9,
    ARetro_m2=0.01,
    attenuation_c=0.05)
H = model.gain(link, env)
print(f"Table XI Retro gain = {H:.6e}")
\end{lstlisting}

% ---- 28. Table XI Reflective NLoS ----
\subsubsection{Table XI Reflective NLoS (Eq.\,91--92)}
\label{sec:uw_reflective}

\paragraph{Registry key.} \texttt{("underwater", "nlos", "reflective\_case")}

\paragraph{Mathematical form.}
The propagation distance through water is $z = h + x$ (transmitter depth $h$
plus horizontal offset $x$).  For incidence angle $\phi$ relative to the
water--air interface:
\begin{equation}
\theta_t = \arcsin\!\left(\frac{n_1}{n_2}\sin\phi\right)
\end{equation}

\textbf{Case 1} ($\phi < \theta_c$, partial transmission via Fresnel):
\begin{equation}
\mathcal{F}(\phi) = \frac{1}{2}\left[
  \left(\frac{\tan(\theta_t - \phi)}{\tan(\theta_t + \phi)}\right)^2
  + \left(\frac{\sin(\phi - \theta_t)}{\sin(\phi + \theta_t)}\right)^2
\right]
\end{equation}

\textbf{Case 2} ($\phi \ge \theta_c$, total internal reflection):
$\mathcal{F} = 1$.

The received power is:
\begin{equation}
\boxed{P_R = P_T\,\eta_T\,\eta_R\;
\frac{A_{\text{rec}}\,\cos\phi}
     {2\pi z^3 (\cos\theta_{\min} - \cos\theta_{\max})}
\; e^{-c\,z} \; \mathcal{F}(\phi)}
\end{equation}

\paragraph{Code snippet.}
\begin{lstlisting}
from owcm.underwater.models.survey_model import TableXIReflective

model = TableXIReflective(link, env,
    h_m=10.0, x_m=10.0,
    phi_rad=np.deg2rad(20),
    theta_min_rad=0.0,
    theta_max_rad=np.deg2rad(45),
    n1=1.33, n2=1.0,
    etaT=0.9, etaR=0.9,
    attenuation_c=0.05)
H = model.gain(link, env)
print(f"Table XI Reflective NLoS gain = {H:.6e}")
\end{lstlisting}


%=============================================================================
%=============================================================================
\newpage
\section{Summary Table}
%=============================================================================
%=============================================================================

\begin{table}[H]
\centering
\caption{Implemented environment models in \texttt{owcm}.
\emph{(D) = deterministic gain only; (S) = stochastic sampling supported.}}
\label{tab:summary}
\footnotesize
\begin{tabular}{@{}llll@{}}
\toprule
\textbf{Environment} & \textbf{Category} & \textbf{Model Key} & \textbf{Type} \\
\midrule
Atmospheric & Attenuation & \texttt{beer\_lambert} & D \\
            & Attenuation & \texttt{kim} & D \\
            & Attenuation & \texttt{kruse} & D \\
            & Attenuation & \texttt{al\_naboulsi\_convection} & D \\
            & Attenuation & \texttt{al\_naboulsi\_advection} & D \\
            & Attenuation & \texttt{ijaz} & D \\
            & Attenuation & \texttt{esmail} & D \\
            & Geometry    & \texttt{uniform\_aperture} & D \\
            & Geometry    & \texttt{gaussian\_collected\_exact} & D \\
            & Geometry    & \texttt{gaussian\_collected\_approx} & D \\
            & Pointing    & \texttt{rayleigh\_gaussian} & S \\
            & Turbulence  & \texttt{lognormal} & S \\
            & Turbulence  & \texttt{gamma\_gamma} & S \\
\midrule
Indoor      & LoS       & \texttt{lambertian\_los} & D \\
            & LoS       & \texttt{lambertian\_los\_plus\_nlos} & D \\
            & NLoS      & \texttt{kbounce} & D \\
            & Blockage  & \texttt{onoff} & S \\
            & Blockage  & \texttt{markov} & S \\
\midrule
Space       & Free Space  & \texttt{fspl} & D \\
            & Deep Space  & \texttt{table\_xiii} & D \\
\midrule
Underground & LoS       & \texttt{lambertian\_los} & D \\
            & Dust      & \texttt{exp\_extinction} & D \\
            & Blockage  & \texttt{onoff} & S \\
            & Blockage  & \texttt{markov} & S \\
\midrule
Underwater  & Extinction & \texttt{a\_plus\_b} / \texttt{c\_direct} & D \\
            & Link       & \texttt{los\_case} (Eq.\,89) & D \\
            & Link       & \texttt{retroreflective\_case} (Eq.\,90) & D \\
            & NLoS       & \texttt{reflective\_case} (Eq.\,91--92) & D \\
\bottomrule
\end{tabular}
\end{table}

\end{document}
